% V84 JEI evidence-frozen manuscript.
% RA-RepDet denotes the sample-dependent dynamic gate. Modality dropout is optional.
% No V84 access to the locked 837-image internal holdout occurred.
\IfFileExists{figures/fig1_method.pdf}{}{\PassOptionsToPackage{demo}{graphicx}}
\documentclass[pdflatex,sn-basic,Numbered,iicol]{sn-jnl}

\usepackage{graphicx}
\usepackage{xcolor}
\usepackage{booktabs}
\usepackage{tabularx}
\usepackage{array}
\usepackage{multirow}
\usepackage{amsmath,amssymb,amsfonts}
\hypersetup{hidelinks}

\newcolumntype{Y}{>{\raggedright\arraybackslash}X}
\newcommand{\figref}[1]{Fig.~\ref{#1}}
\newcommand{\tabref}[1]{Table~\ref{#1}}
\newcommand{\ap}{\mathrm{AP}}

\title[RA-RepDet for tri-modal UAV detection]{RA-RepDet: Reliability-Aware RGB-Thermal-Event Fusion with Leakage-Aware Validation and Supervised Cross-Dataset Transfer}

\author[1]{\fnm{Nan} \sur{Xin}}
\author*[1]{\fnm{Xueting} \sur{Jin}}\email{jinxueting@ahpc.edu.cn}
\affil*[1]{\orgname{Anhui Police College}, \orgaddress{\city{Hefei}, \state{Anhui}, \country{China}}}

\abstract{Multimodal UAV detection can benefit from RGB, thermal, and event inputs, but related frames can leak across splits and robustness gains can conflate adaptive fusion with modality dropout. We present RA-RepDet, a lightweight RepViT-FPN-FCOS detector whose defining mechanism is sample-dependent dynamic gating over modality-specific features. On a 2,213-image component-disjoint TriAir development-validation split, gate/no-dropout reaches $0.7251\pm0.0121$ COCO AP, exceeding matched early fusion by $0.0449\pm0.0341$. A separately trained RGB+thermal baseline, under identical optimization, selection, and evaluation contracts, reaches $0.6843\pm0.0312$ AP, so this experiment does not isolate a positive incremental event contribution. In a matched $2\times2$ analysis, gating improves full-input AP by 0.0382 on average, whereas robustness after event removal is mainly associated with modality-dropout training. Component-cluster bootstrap over 1,298 leakage components yields positive 95\% intervals for gate/no-dropout against early, fixed-equal, and deterministic-projection controls. Controlled blur and attenuation reduce AP but do not produce monotonic decreases in the affected modality's gate weight; weights are therefore task-driven fusion coefficients, not calibrated sensor-health estimates. Supervised MM-UAV transfer remains descriptive exposed-devval evidence. These results support dynamic fusion under leakage-aware validation, not physical reliability, statistical significance, or independent external generalization.}

\keywords{UAV vehicle detection; RGB-thermal-event fusion; reliability-aware fusion; modality dropout; component-disjoint validation; supervised cross-dataset transfer; multi-sensor perception.}

\begin{document}
\raggedbottom
\maketitle

\section{Introduction}
UAV vehicle detection must localize small targets under changing altitude, viewpoint, scale, illumination, motion, and clutter. These conditions are challenging for generic object detection and for aerial imaging in particular \cite{zhao2019object,cheng2016survey,xia2022dota,zhu2017deepRS}. RGB imagery provides texture and color, but may weaken under poor contrast or changing exposure. Thermal observations can provide temperature-related contrast, while event representations can encode brightness changes with high temporal resolution and wide dynamic range \cite{baltrusaitis2019multimodal,gallego2022event}. The key design question is therefore not merely how to concatenate modalities, but how to use their complementary information without unnecessarily increasing the cost of a lightweight detector.

A second issue is evaluation integrity. In multimodal UAV recordings, nearby observations can remain visually similar even when their file contents differ. If such observations occur on opposite sides of a split, validation performance can be optimistic. This risk is distinct from model design and should be addressed before interpreting a fusion gain \cite{kapoor2023leakage,cawley2010overfitting}. In the present project, split review first identified exact decoded-RGB duplicates and then identified additional adjacent-or-near-identical observations through a human review workflow. These findings motivate a more conservative validation protocol, not a claim that every possible dependence has been eliminated.

This paper presents RA-RepDet, a compact five-channel RGB-thermal-event detector. RA-RepDet is defined by separate modality stems and an image-conditioned softmax gate feeding a RepViT-M0.9 \cite{wang2024repvit}, feature-pyramid network (FPN) \cite{lin2017fpn}, and FCOS detector \cite{tian2019fcos}. Gate/no-dropout is the primary nominal-accuracy variant. Independent modality dropout at $p=0.15$ is evaluated separately as an optional robustness regularizer; it is not part of the definition of RA-RepDet and is not selected as the nominal default.

The study first evaluates the fusion design under a matched TriAir sensing configuration and a component-disjoint validation protocol. It then examines whether the resulting source-domain representation transfers to a different tri-modal acquisition system. The MM-UAV study deliberately separates an unadapted frozen-model stress test from a supervised target-domain benchmark with learned feature alignment. This separation is important because cross-sensor geometry, modality representation, and optimization effects can otherwise be conflated with the value of the fusion mechanism itself.

The conclusions are intentionally scoped. The TriAir evidence is validation-only, synthetic channel removal is not a physical sensor-failure experiment, the frozen MM-UAV experiment uses an unregistered adapter, and the supervised MM-UAV benchmark uses target-domain labels on an exposed devval split. None of these experiments is presented as independent untouched-test evidence.

\noindent\textbf{Contributions.}
\begin{itemize}
    \item We define RA-RepDet as a lightweight sample-dependent dynamic gate over modality-specific features with a shared RepViT-FPN-FCOS detector.
    \item We report three-seed component-disjoint validation, a separately trained RGB+thermal control, and a matched $2\times2$ gate-by-dropout channel-removal analysis.
    \item We analyze gate weights under clean inputs and controlled modality degradation, and report component-cluster bootstrap intervals for three primary dynamic-gate comparisons.
    \item We retain checkpoint-backed single-modality and supervised MM-UAV evidence with explicit development-validation, alignment, transfer, and dissemination boundaries.
\end{itemize}

\section{Related Work}
\subsection{Aerial Detection and Lightweight Detectors}
Aerial detection involves wide scale variation, small objects, geometric changes, and cluttered backgrounds \cite{cheng2016survey,zhu2017deepRS,xia2022dota}. Modern detectors include region-proposal and dense one-stage formulations \cite{ren2017faster,lin2020focal}. FCOS is an anchor-free dense detector that predicts object locations directly \cite{tian2019fcos}. In this study, the detector head is held fixed; the research question is whether the tri-modal fusion front end can improve a compact detection pipeline. RepViT is used as a mobile-oriented backbone, and FPN provides multi-scale features \cite{wang2024repvit,lin2017fpn}.

\subsection{Multimodal and Event-Based Perception}
Multimodal learning integrates complementary sources while accounting for differences in their reliability and representation \cite{baltrusaitis2019multimodal,ramachandram2017deep}. Visible-infrared fusion has been studied extensively at pixel and feature levels \cite{li2017pixel,ma2019surveyfusion,li2019densefuse}. Event cameras provide asynchronous brightness-change information and can support perception under fast motion or high dynamic range \cite{lichtsteiner2008sensor,gallego2022event,gehrig2021dsec}. RA-RepDet does not reconstruct raw event streams or infer a missing modality. It consumes a preconstructed five-channel sample and learns a compact image-level mixture of modality-stem features.

\subsection{Missing Modality and Evaluation Validity}
Modality dropout is a common strategy for limiting brittle dependence on one information source \cite{neverova2016moddrop}. Here, it is used as a training-time intervention, while zeroing one channel group at evaluation time gives a controlled synthetic channel-removal protocol. This protocol is useful for comparing models under the same intervention, but it cannot reproduce sensor drift, misregistration, field-of-view change, or other physical failure modes.

Evaluation leakage and adaptive reuse of validation information can produce optimistic performance estimates \cite{kapoor2023leakage,cawley2010overfitting}. The present work does not claim that any finite audit proves the absence of all dependence. Instead, it constructs a split whose train-validation boundary does not cut components of a specified graph formed from exact RGB, pHash, dHash, and human-adjudicated adjacent-or-near-identical relations.

\subsection{Transfer Learning and Cross-Domain Adaptation}
Transfer learning reuses information learned from a source task or domain to support a target task, but its benefit depends on the relationship between the source and target distributions \cite{pan2010transfer}. Deep features tend to become more task-specific in later layers, and larger source-target differences can reduce their transferability \cite{yosinski2014transfer}; source knowledge can therefore help or have limited effect under different target protocols. Multimodal cross-sensor transfer adds a further complication because modalities may differ in appearance statistics, spatial calibration, field of view, event construction, and annotation geometry. Our protocol consequently evaluates source initialization only after fixing a target-domain alignment-aware architecture and compares it with a matched from-scratch control.

\section{Method}
\subsection{Problem Formulation and Detector Interface}
Each observation is a five-channel tensor
\begin{equation}
\mathbf{x}=[\mathbf{x}_{\mathrm{rgb}},\mathbf{x}_{\mathrm{th}},\mathbf{x}_{\mathrm{ev}}],
\end{equation}
where $\mathbf{x}_{\mathrm{rgb}}\in\mathbb{R}^{3\times H\times W}$, $\mathbf{x}_{\mathrm{th}}\in\mathbb{R}^{1\times H\times W}$, and $\mathbf{x}_{\mathrm{ev}}\in\mathbb{R}^{1\times H\times W}$. The task is single-class vehicle detection. Raw TriAir class 0 is mapped to the foreground label used by the detector, with background retained as label 0.

The downstream detector is shared by all TriAir fusion controls: a RepViT-M0.9 backbone, an FPN with 128 output channels, and an FCOS head. This shared interface supports matched comparisons among early fusion, fixed or deterministic stem fusion, and the sample-dependent RA-RepDet gate.

\begin{figure*}[t]
\centering
\includegraphics[width=0.98\textwidth]{figures/fig1_method.pdf}
\caption{RA-RepDet. RGB, thermal, and event inputs are first encoded by modality-specific stems. A small image-level reliability module produces softmax fusion weights, the weighted feature is projected to the shared detector interface, and a RepViT-FPN-FCOS stack produces vehicle detections. Modality dropout is used only during training.}
\label{fig:method}
\end{figure*}

\subsection{Matched Early Fusion}
The matched early-fusion system projects the five input channels to three channels with a learned $1\times1$ convolution before the shared RepViT-FPN-FCOS detector. It has access to all five channels but does not form explicit modality-specific features or sample-dependent fusion weights.

\subsection{Reliability-Aware Fusion}
The reliability-aware model processes RGB, thermal, and event inputs through separate $3\times3$ convolution--batch-normalization--SiLU stems. Each stem produces a 16-channel feature map, denoted by $\mathbf{F}_{\mathrm{rgb}}$, $\mathbf{F}_{\mathrm{th}}$, and $\mathbf{F}_{\mathrm{ev}}$. Global average pooling yields one descriptor per stem. The descriptors are concatenated and passed through a small two-layer multilayer perceptron to generate three logits. The fusion weights are
\begin{equation}
\begin{aligned}
\boldsymbol{\alpha}=\operatorname{softmax}\bigl(g([&\operatorname{GAP}(\mathbf{F}_{\mathrm{rgb}});\\
&\operatorname{GAP}(\mathbf{F}_{\mathrm{th}});\operatorname{GAP}(\mathbf{F}_{\mathrm{ev}})])\bigr),
\end{aligned}
\end{equation}
where $\boldsymbol{\alpha}=[\alpha_{\mathrm{rgb}},\alpha_{\mathrm{th}},\alpha_{\mathrm{ev}}]$ and the three entries sum to one. The fused feature is
\begin{equation}
\mathbf{F}_{\mathrm{fuse}}=\alpha_{\mathrm{rgb}}\mathbf{F}_{\mathrm{rgb}}+\alpha_{\mathrm{th}}\mathbf{F}_{\mathrm{th}}+\alpha_{\mathrm{ev}}\mathbf{F}_{\mathrm{ev}}.
\end{equation}
A projection converts $\mathbf{F}_{\mathrm{fuse}}$ to the three-channel interface expected by the shared backbone. The gate is a learned feature-mixture mechanism; its weights are not calibrated probabilities of physical sensor health.

\subsection{Optional Modality-Dropout Training}
RGB channels $0{:}3$, thermal channel 3, and event channel 4 can be independently zeroed during training with probability $p=0.15$. If all three groups are selected, one group is restored uniformly at random so that at least one stream remains. This intervention is evaluated in both early and gated architectures to separate its effect from dynamic gating. Ordinary all-modal inference supplies all streams, and gate/no-dropout is the primary RA-RepDet nominal-accuracy configuration.

\section{Dataset and Component-Disjoint Validation Protocol}
\subsection{TriAir Provider Archive and Runtime Representation}
The dataset paper introduces 10,489 synchronized and pre-aligned RGB--thermal--event frames and reports 24,223 annotated vehicles \cite{iaboni2026trimodal}. For this study, the local tree at \path{D:\download\triair} was compared entry by entry with the provider's current \path{triair.zip}: all 20,240 archive entries matched in relative path, uncompressed size, and CRC32, with zero missing, extra, or different entries. The provider assigns no semantic version or release tag, so the study records the source as the \emph{untagged provider archive}, Google Drive file ID \path{1w71v6n41yqjP7BCr9ni4JdcxMnQ2ocR0}, Last-Modified 2025-11-21, archive size 3,551,150,083 bytes. These archive-identity facts are from the authors' complete local audit; the external provider page exposes the same file identity but does not publish a semantic version.

The provider archive itself contains 10,489 five-channel \path{.npy} arrays and 9,751 YOLO \path{.txt} files; no project-side offline conversion from raw RGB, thermal, or event recordings was found. Every audited array has shape $(301,391,5)$ and data type \path{uint8}, with RGB in channels 0--2, thermal in channel 3, and event representation in channel 4, consistent with the provider paper and upstream repository. At runtime only, the project converts HWC to CHW, casts \path{uint8} values to floating point and divides by 255, converts normalized YOLO boxes to absolute \path{xyxy}, and maps raw class 0 to detector foreground class 1. Missing and empty label files remain valid empty-target observations.
